\section{Interrogare il database}
Abbiamo parlato di come descriviamo e come conserviamo 
questi famigerati \textit{dati}, ma come lavoriamo su quest'ultimi? \\
Useremo un linguaggio interattivo: \textbf{SQL}. Quest'ultimo
è un linguaggio standard oramai per interfacciarsi con database e
lo troveremo già incluso in tanti linguaggi come \textbf{C} 
oppure \textbf{Pascal}.
\\
\\
Il linguaggio è composto principalmente da due tipologie di operazioni:
\textbf{DML} e \textbf{DDL}. \\
Le operazioni di \textbf{Data Manipulation Language} manipolano 
le istanze del databse, e quindi modificano i dati al suo interno;
il secondo tipo di operazioni, \textbf{Data Definition Language}, 
definiscono lo schema del database e come i dati si relazionino 
tra di loro. \\
Seguono degli esempi di, rispettivamente, operazione di DDL e DML

\begin{verbatim}
  CREATE TABLE orario (
    insegnamento CHAR(20),
    docente CHAR(20),
    aula CHAR(4),
    ora CHAR(5)
  );
\end{verbatim}

\begin{verbatim}
  INSERT INTO orario (insegnamento, docente, aula, ora)
  VALUES ('Basi di dati', 'Mario Rossi', 'N1', '09:00');
\end{verbatim}

Questi esempi servono soltanto per avere un rapido chiarimento sulle
differneze tra i due tipi di comandi, ovviamente una persona che non ha
mai usato SQL avrà difficoltà a capire. Basta sapere che il primo comando
\verb|CREATE TABLE| definisce lo schema della tabella \textit{orario},
il secondo comando (\verb|INSERT INTO|), va invece a modificare questa 
tabella aggiungendo dei nuovi dati.